% Seminar 4 handout (lecture 4): smooth numbers, relations, the linear system modulo p - 1,
% the individual logarithm.
\documentclass{ca-handout}

\seminar{4}{Метод исчисления индексов}

\begin{document}
\makehandouttitle

\section{Гладкие числа и факторная база}
Число $m$ \emph{$B$-гладкое}, если все его простые делители $\le B$; $\psi(x, B)$ "--- число
$B$-гладких $m \le x$. При $u = \ln x / \ln B$ доля гладких чисел $\psi(x, B)/x \approx u^{-u}$
(Кэнфилд"--~Эрдёш"--~Померанс). \emph{Факторная база} $F = \{\ell \le B : \ell \text{ простое}\}$,
$t = \pi(B)$; неизвестные $L_\ell = \log_g \ell \in \Zn{(p - 1)}$, $g$ "--- первообразный корень.

\begin{reminder}[title={Соотношение и логарифм элемента}]
\[
\begin{aligned}
  g^k \bmod p &= \prod_{\ell \in F} \ell^{e_\ell} &&\Longrightarrow&
  \sum_{\ell} e_\ell L_\ell &\equiv k \pmod{p - 1},\\
  h g^{s} \bmod p &= \prod_{\ell \in F} \ell^{f_\ell} &&\Longrightarrow&
  \log_g h &\equiv \sum_{\ell} f_\ell L_\ell - s \pmod{p - 1}.
\end{aligned}
\]
Значения $g^k$ (и $h g^s$) удобно получать последовательно: следующее "--- предыдущее, умноженное
на $g$. В варианте $h g^{-s}$: $\log_g h \equiv s + \sum f_\ell L_\ell$.
\end{reminder}

\begin{reminder}[title={Факторная база с $-1$}]
Раскладывают абсолютно наименьший вычет $r \equiv g^k$, $\abs{r} < p/2$; знак даёт слагаемое
$L_{-1}$. Так как $g^{(p-1)/2}$ "--- элемент порядка $2$, т.\,е. $-1$, то
\[
  L_{-1} = \log_g(-1) = \frac{p - 1}{2}.
\]
\end{reminder}

\section{Система по модулю \texorpdfstring{$p - 1$}{p - 1}}
$\Zn{(p-1)}$ "--- не поле. Систему решают по модулю каждого $q^e \parallel p - 1$ и собирают
по КТО: $x \equiv \sum_i r_i M_i y_i \pmod M$, $M_i = M/m_i$, $y_i \equiv M_i^{-1} \pmod{m_i}$.

\begin{reminder}[title={Метод Гаусса по модулю $q^e$}]
\begin{itemize}
  \item Ведущий элемент должен быть \emph{обратим}, т.\,е. не делиться на $q$; строку умножают на
    обратный к нему и вычитают из остальных строк.
  \item Решение единственно по модулю $q^e$ $\iff$ ранг матрицы по модулю $q$ равен $t$ $\iff$
    в каждом столбце находится обратимый ведущий элемент. Для квадратной системы: $q \nmid \det$.
  \item Сравнение $c L_\ell \equiv k \pmod{q^e}$ с $q \mid c$ определяет $L_\ell$ лишь по модулю
    $q^{e - v_q(c)}$ (если вообще совместно); нужны другие соотношения.
  \item Обычно хватает $t + s$ соотношений с небольшим запасом $s$; чаще всего система
    вырождена по модулю $2$.
\end{itemize}
\end{reminder}

\begin{reminder}[title={Проверка и подсказка по модулю малых $q$}]
$L_\ell$ чётен $\iff$ $\legendre{\ell}{p} = 1$ (критерий Эйлера: $\ell^{(p-1)/2} = (-1)^{L_\ell}$).
Вообще, $L_\ell \bmod q^e$ "--- логарифм $\ell^{(p-1)/q^e}$ по основанию $g^{(p-1)/q^e}$
(Полиг"--~Хеллман). Напоминание: $\legendre{2}{p} = 1 \iff p \equiv \pm 1 \pmod 8$;
$\legendre{q}{p} = \legendre{p}{q}$, если $p \equiv 1$ или $q \equiv 1 \pmod 4$
($\legendre{q}{p} = -\legendre{p}{q}$, если оба $\equiv 3 \pmod 4$).
\end{reminder}

\section{Алгоритм и выбор \texorpdfstring{$B$}{B}}
\begin{enumerate}
  \item Выбрать $B$, факторную базу $F$.
  \item Перебирая $k$, собрать $\approx t + s$ соотношений $g^k = \prod \ell^{e_\ell}$.
  \item Решить систему по модулю каждого $q^e \parallel p - 1$, собрать $L_\ell$ по КТО;
    проверить $g^{L_\ell} \equiv \ell$.
  \item Найти $s$ с гладким $h g^{s}$ и вычислить $\log_g h$.
\end{enumerate}

\begin{reminder}[title={$L$-нотация и сложность (эвристически)}]
$L_p[\alpha, c] = \exp\bigl((c + o(1))(\ln p)^{\alpha}(\ln\ln p)^{1-\alpha}\bigr)$,
$L(p) = e^{\sqrt{\ln p \ln\ln p}}$. При $B = L_p[1/2, c]$ доля гладких $\approx
L_p[1/2, -\frac{1}{2c}]$, попыток на этапе 1: $L_p[1/2, c + \frac{1}{2c}]$. Оптимум:
\[
  B \approx L(p)^{1/\sqrt 2}, \qquad \text{этап 1: } L_p[1/2, \sqrt 2], \qquad
  \text{этап 2: } L_p[1/2, 1/\sqrt 2].
\]
Для малых $p$ рецепт лишь ориентир: $B$ округляют, оставляя небольшую систему.
Решето числового поля: $L_p[1/3, 1{,}923]$. Для эллиптических кривых общего вида метод
неприменим.
\end{reminder}
\end{document}
